% Capitulo: O data warehouse e a modelagem dimensional
\chapter{O data warehouse e a modelagem dimensional}
\label{cap:dw}

Nos anos 1990, o BI deixou de ser um conjunto de relatórios e planilhas e virou projeto
corporativo, com orçamento, equipe e arquitetura próprios. O centro desse projeto era o
data warehouse, o armazém de dados. Este capítulo conta como ele surgiu, apresenta as
duas escolas que definiram o campo, a de Bill Inmon e a de Ralph Kimball, e mostra por
que a modelagem dimensional, uma ideia dessa época, continua sendo uma das habilidades
mais úteis para quem trabalha com dados.

\section{Como surgiu o data warehouse}
\label{sec:surgimento-dw}

Os sistemas operacionais de uma empresa (vendas, estoque, contabilidade, folha de
pagamento) foram construídos para registrar transações. Eles precisam gravar um pedido
em milissegundos, garantir que dois vendedores não vendam a mesma unidade e manter o
saldo sempre correto. Para isso, guardam o estado atual das coisas, organizado de forma
a evitar repetição.

Análise pede quase o oposto. O gestor quer saber como as vendas evoluíram nos últimos
três anos, por região e por produto, comparando com o orçamento. Essa consulta lê
milhões de linhas, precisa de histórico que o sistema operacional muitas vezes já
apagou, e cruza dados de sistemas que não conversam entre si. Rodar esse tipo de
consulta no mesmo banco que registra as vendas deixa a operação lenta, e nem assim
resolve o cruzamento entre sistemas.

A saída que as empresas encontraram primeiro foi extrair dados. Cada área pedia um
arquivo do sistema de origem, tratava do seu jeito e montava o seu relatório. Com o
tempo, as extrações se multiplicavam, cada uma com uma regra um pouco diferente, e os
números de áreas diferentes deixavam de bater. \pendente{achar fonte para a
proliferação de extrações; Inmon descreve esse cenário}

Em 1988, Barry Devlin e Paul Murphy descreveram na IBM uma arquitetura diferente: um
armazém de dados de negócio, separado dos sistemas operacionais, alimentado de forma
controlada e usado como fonte única para os relatórios \parencite{devlin1988}. Quatro
anos depois, Bill Inmon publicou o livro que popularizou o termo data warehouse e
transformou a ideia numa disciplina \parencite{inmon1992}.

A definição de Inmon ficou clássica. Um data warehouse é uma coleção de dados orientada
por assunto, integrada, variável no tempo e não volátil, que apoia a tomada de decisão
gerencial \parencite{inmon2005}. \pendente{conferir a citação exata e a página na
edição usada} Cada um dos quatro adjetivos resolve um problema concreto:

\begin{description}[style=nextline, leftmargin=1.5em, font=\normalfont\itshape]
  \item[Orientada por assunto] Os dados se organizam pelos assuntos do negócio
    (cliente, produto, venda), não pelos sistemas de onde vieram. O analista procura
    \enquote{clientes} e não precisa saber se eles estão no ERP ou no CRM.
  \item[Integrada] O mesmo cliente tem o mesmo código e o mesmo nome em todo o armazém,
    mesmo que cada sistema de origem o registre de um jeito. É a primeira aparição, no
    livro, do problema de master data.
  \item[Variável no tempo] O armazém guarda o histórico. Se o cliente mudou de cidade, é
    possível saber onde ele estava quando fez cada compra.
  \item[Não volátil] O dado carregado não é alterado pelas transações do dia a dia. Uma
    consulta de hoje sobre o mês passado dá o mesmo resultado amanhã.
\end{description}

A Figura~\ref{fig:dw-linha-tempo} situa esses marcos numa linha do tempo que vai até a
IA. O resto do livro percorre, de certa forma, a metade de baixo dessa linha.

\begin{figure}[htbp]
  \centering
  % Linha do tempo do data warehouse (cap. 2). Fontes na legenda do capitulo.
\begin{tikzpicture}[diagrama, x=1mm, y=1mm]
  \draw[fio, line width=0.8pt] (18,5) -- (18,-85);
  \foreach \y/\ano/\texto/\peso in {
      0/1958/{Luhn descreve um \enquote{sistema de business intelligence} na IBM}/normal,
    -10/1970/{Codd propõe o modelo relacional, base dos bancos de dados SQL}/normal,
    -20/1988/{Devlin e Murphy descrevem a arquitetura de um \enquote{business data warehouse} na IBM}/normal,
    -30/1992/{Inmon publica \emph{Building the Data Warehouse}: o armazém corporativo, normalizado, construído de cima para baixo}/destaque,
    -40/1996/{Kimball publica \emph{The Data Warehouse Toolkit}: a modelagem dimensional, construída de baixo para cima}/destaque,
    -50/{Anos 2000}/{O Data Vault, de Dan Linstedt, ganha espaço como alternativa para a camada de integração}/normal,
    -60/{Anos 2010}/{Armazéns de dados na nuvem separam armazenamento de processamento, e o ETL dá lugar ao ELT}/normal,
    -70/2021/{Armbrust e colegas propõem o lakehouse, que junta data lake e data warehouse}/normal,
    -80/{Anos 2020}/{A IA generativa entra no fluxo de trabalho com dados}/normal%
  }{
    \ifthenelse{\equal{\peso}{destaque}}
      {\filldraw[draw=tinta, fill=tinta] (18,\y) circle (1.4);}
      {\filldraw[draw=cinza, fill=white, line width=0.6pt] (18,\y) circle (1.1);}
    \node[anchor=east, font=\footnotesize] at (14.5,\y) {\ano};
    \node[anchor=west, text width=92mm, align=left, font=\footnotesize] at (22,\y) {\texto};
  }
\end{tikzpicture}

  \caption{Uma linha do tempo do data warehouse.}
  \label{fig:dw-linha-tempo}
  \fonte{elaborado pelo autor com base em \textcite{luhn1958}, \textcite{codd1970},
    \textcite{devlin1988}, \textcite{inmon1992}, \textcite{kimball1996},
    \textcite{linstedt2015} e \textcite{armbrust2021}.}
\end{figure}

\section{A abordagem de Inmon}
\label{sec:inmon}

Para Inmon, o data warehouse é a base corporativa de informação. Ele deve ser
construído de cima para baixo: primeiro se modela a empresa como um todo, depois se
carrega um armazém central, integrado e normalizado, na terceira forma normal, o mesmo
tipo de modelagem usado nos sistemas transacionais para evitar repetição de dados
\parencite{inmon2005}. Esse armazém é a fonte de verdade.

Os data marts, bases menores voltadas a uma área (vendas, finanças, logística), vêm
depois e são derivados do armazém central. Por isso são chamados de dependentes. Neles,
sim, os dados podem ser reorganizados para facilitar a consulta. A
Figura~\ref{fig:inmon} resume o desenho.

\begin{figure}[htbp]
  \centering
  % Arquitetura de Inmon: de cima para baixo (cap. 2).
\begin{tikzpicture}[diagrama, x=1mm, y=1mm]
  % sistemas de origem
  \foreach \y/\nome/\id in {13.5/ERP/f1, 4.5/CRM/f2, -4.5/Planilhas/f3, -13.5/{Outros}/f4}
    \node[caixa, minimum width=15mm, minimum height=6mm] (\id) at (8,\y) {\nome};
  \node[moldura, fit=(f1)(f4), label={[rotulo]above:Sistemas de origem}] {};

  % armazem corporativo
  \node[bancodestaque, minimum width=26mm, minimum height=26mm] (edw) at (46,0)
    {Data warehouse\\corporativo\\[1mm]{\scriptsize normalizado (3FN)}\\{\scriptsize e integrado}};
  \foreach \id in {f1,f2,f3,f4} \draw[seta] (\id.east) -- (edw.west);
  \node[rotulo] at (25.5,18) {ETL};

  % data marts dependentes
  \foreach \y/\nome/\id in {13/Vendas/m1, 0/Finanças/m2, -13/Logística/m3}
    \node[banco, minimum width=17mm, minimum height=9mm] (\id) at (79,\y) {\nome};
  \node[moldura, fit=(m1)(m3), label={[rotulo]above:Data marts dependentes}] {};
  \foreach \id in {m1,m2,m3} \draw[seta] (edw.east) -- (\id.west);

  % consumo
  \foreach \y/\nome/\id in {13/Relatórios/u1, 0/Dashboards/u2, -13/Análises/u3}
    \node[caixa, minimum width=16mm, minimum height=6mm] (\id) at (108,\y) {\nome};
  \foreach \i in {1,2,3} \draw[seta] (m\i.east) -- (u\i.west);

  \node[rotulo, text width=110mm] at (58,-27)
    {Ordem de construção: primeiro o armazém corporativo, depois os data marts de cada área.};
\end{tikzpicture}

  \caption{A arquitetura de Inmon: armazém corporativo normalizado e data marts dependentes.}
  \label{fig:inmon}
  \fonte{elaborado pelo autor com base em \textcite{inmon2005} e \textcite{inmon2001}.}
\end{figure}

Inmon e coautores ampliaram depois essa visão para a Corporate Information Factory, uma
arquitetura completa em que o data warehouse é uma das peças, ao lado de bases
operacionais integradas, data marts e aplicações de análise \parencite{inmon2001}.

A força dessa abordagem é a consistência. Como tudo passa pelo armazém central, as áreas
partem dos mesmos dados e das mesmas regras. E como o núcleo é normalizado, ele aguenta
bem perguntas que ninguém previu quando o modelo foi desenhado. O ponto fraco é o tempo.
Modelar a empresa inteira antes de entregar o primeiro relatório exige meses, às vezes
anos, e um patrocínio que nem sempre sobrevive à espera.

\section{A abordagem de Kimball}
\label{sec:kimball}

Ralph Kimball propôs o caminho inverso. Em vez de começar pela empresa inteira, o
armazém é construído um processo de negócio por vez: primeiro vendas, depois estoque,
depois compras. Cada processo ganha um modelo dimensional, desenhado para ser fácil de
entender e rápido de consultar \parencite{kimball1996,kimball2013}.

O que impede esse crescimento incremental de virar um amontoado de bases isoladas são as
dimensões conformadas. Data, produto, cliente e loja são as mesmas tabelas, com as mesmas
chaves e os mesmos atributos, em todos os modelos. Assim, uma análise que cruza vendas e
estoque por produto funciona, porque os dois usam exatamente a mesma dimensão de produto.
O conjunto de modelos ligados pelas dimensões conformadas forma o que Kimball chama de
arquitetura de barramento (Figura~\ref{fig:kimball}). Uma dimensão conformada é, na
prática, master data com outro nome, e o capítulo~\ref{cap:master-data} volta a esse
ponto.

\begin{figure}[htbp]
  \centering
  % Arquitetura de Kimball: data marts dimensionais integrados pelo barramento (cap. 2).
\begin{tikzpicture}[diagrama, x=1mm, y=1mm]
  % sistemas de origem
  \foreach \y/\nome/\id in {13.5/ERP/f1, 4.5/CRM/f2, -4.5/Planilhas/f3, -13.5/{Outros}/f4}
    \node[caixa, minimum width=15mm, minimum height=6mm] (\id) at (8,\y) {\nome};
  \node[moldura, fit=(f1)(f4), label={[rotulo]above:Sistemas de origem}] {};

  % area de preparacao
  \node[caixa, minimum width=19mm, minimum height=30mm] (stg) at (36,0)
    {Área de\\preparação\\[1mm]{\scriptsize (ETL)}};
  \foreach \id in {f1,f2,f3,f4} \draw[seta] (\id.east) -- (\id.east -| stg.west);

  % area de apresentacao: data marts dimensionais
  \foreach \x/\nome/\id in {62/Vendas/m1, 80/Estoque/m2, 98/Finanças/m3}
    \node[bancodestaque, minimum width=16mm, minimum height=13mm] (\id) at (\x,5)
      {\nome\\{\scriptsize estrela}};
  % barramento de dimensoes conformadas
  \node[caixa, fill=fundo, draw=cinza, minimum width=56mm, minimum height=6mm, inner sep=1pt] (bus) at (80,-12)
    {\scriptsize Barramento: dimensões conformadas (Data, Produto, Cliente)};
  \foreach \id in {m1,m2,m3} \draw[ligacao, line width=1.2pt] (\id.south) -- (\id.south |- bus.north);
  \node[moldura, fit=(m1)(m3)(bus), label={[rotulo]above:Área de apresentação}] (apr) {};
  \draw[seta] (stg.east) -- (stg.east -| apr.west);

  \node[rotulo, text width=110mm] at (58,-27)
    {Ordem de construção: um processo de negócio por vez, integrados pelas dimensões conformadas.};
\end{tikzpicture}

  \caption{A arquitetura de Kimball: data marts dimensionais integrados pelo barramento de dimensões conformadas.}
  \label{fig:kimball}
  \fonte{elaborado pelo autor com base em \textcite{kimball2013}.}
\end{figure}

O planejamento desse barramento cabe numa tabela simples, a matriz de barramento: os
processos de negócio nas linhas, as dimensões nas colunas, e uma marca onde o processo
usa a dimensão (Tabela~\ref{tab:matriz-barramento}). Ela mostra, de relance, quais
dimensões precisam ser construídas com mais cuidado, porque serão compartilhadas por
muitos processos.

\begin{table}[htbp]
  \centering
  \caption{Uma matriz de barramento: processos de negócio nas linhas, dimensões conformadas nas colunas.}
  \label{tab:matriz-barramento}
  \small
  \begin{tabular}{@{}lccccc@{}}
    \toprule
    Processo de negócio & Data & Produto & Cliente & Loja & Fornecedor \\
    \midrule
    Vendas      & $\bullet$ & $\bullet$ & $\bullet$ & $\bullet$ &           \\
    Compras     & $\bullet$ & $\bullet$ &           & $\bullet$ & $\bullet$ \\
    Estoque     & $\bullet$ & $\bullet$ &           & $\bullet$ &           \\
    Devoluções  & $\bullet$ & $\bullet$ & $\bullet$ & $\bullet$ &           \\
    \bottomrule
  \end{tabular}
  \fonte{exemplo do autor, no formato da matriz de barramento de \textcite{kimball2013}.}
\end{table}

Duas outras ideias de Kimball aparecem em quase todo projeto de dados até hoje. As
chaves substitutas (\emph{surrogate keys}) são códigos gerados pelo próprio armazém para
cada linha de dimensão, independentes dos códigos dos sistemas de origem. Elas protegem o
modelo quando a origem muda a numeração ou quando dois sistemas usam o mesmo código para
coisas diferentes. As dimensões que mudam devagar (\emph{slowly changing dimensions})
são as técnicas para lidar com atributos que mudam com o tempo, como a cidade de um
cliente. A forma mais simples sobrescreve o valor antigo; a mais usada para análise
histórica cria uma nova linha a cada mudança, com datas de início e fim de validade.
Essa segunda forma vai reaparecer, com outro nome, quando o livro tratar de dados para
modelos preditivos (\veja{sec:vazamento}).

A força da abordagem de Kimball é a entrega rápida e o modelo que o usuário entende. O
ponto fraco é que tudo depende da disciplina nas dimensões conformadas. Sem ela, cada
processo ganha a sua própria dimensão de cliente, e o armazém vira um conjunto de silos
com a aparência de um sistema integrado.

\section{Fatos, dimensões e granularidade}
\label{sec:modelagem-dimensional}

O modelo dimensional tem duas peças. As tabelas de fatos guardam as medidas de um
processo de negócio: quantidade vendida, valor, custo. As tabelas de dimensão guardam o
contexto dessas medidas: quando, o quê, quem, onde. Desenhado, o modelo parece uma
estrela, com o fato no centro e as dimensões em volta (Figura~\ref{fig:estrela}).

\begin{figure}[htbp]
  \centering
  % Esquema estrela de vendas (cap. 2). Os mesmos nomes da consulta SQL do capitulo.
\begin{tikzpicture}[diagrama, x=1mm, y=1mm]
  \node[tabela, line width=0.9pt] (f) at (0,0)
    {\textbf{fato\_vendas}\nodepart{second}id\_data\\id\_produto\\id\_cliente\\id\_loja\\[0.6mm]quantidade\\valor\_total};

  \node[tabela] (d1) at (-40,20)
    {\textbf{dim\_data}\nodepart{second}id\_data\\data\\mes\\trimestre\\ano};
  \node[tabela] (d2) at (40,20)
    {\textbf{dim\_produto}\nodepart{second}id\_produto\\nome\\categoria\\marca};
  \node[tabela] (d3) at (-40,-20)
    {\textbf{dim\_cliente}\nodepart{second}id\_cliente\\nome\\cidade\\segmento};
  \node[tabela] (d4) at (40,-20)
    {\textbf{dim\_loja}\nodepart{second}id\_loja\\nome\\cidade\\regiao};

  \foreach \d in {d1,d2,d3,d4} \draw[ligacao, line width=0.7pt] (f) -- (\d);
  \node[rotulo, anchor=north] at (f.south) [yshift=-1mm] {grão: uma linha por item vendido};
\end{tikzpicture}

  \caption{Um esquema estrela de vendas. A tabela de fatos, no centro, guarda as chaves das dimensões e as medidas (quantidade e valor); as dimensões, em volta, guardam o contexto.}
  \label{fig:estrela}
  \fonte{elaborado pelo autor.}
\end{figure}

\textcite{kimball2013} organizam o desenho em quatro passos: escolher o processo de
negócio, declarar o grão, identificar as dimensões e identificar os fatos. A ordem
importa, e o segundo passo é o mais importante e o mais esquecido.

O grão é o que uma linha da tabela de fatos representa. \enquote{Uma linha por item
vendido} é um grão; \enquote{uma linha por pedido} é outro; \enquote{uma linha por loja
por dia} é um terceiro. Cada um responde a perguntas diferentes, e misturar grãos na
mesma tabela é uma das origens mais comuns de número errado em dashboard. Se parte das
linhas é por item e parte é por pedido, qualquer soma conta alguma coisa duas vezes.
Kimball recomenda declarar o grão mais detalhado que os dados permitem, porque agregar
depois é fácil e desagregar é impossível.

\rascunho{
  Relato: um erro de grão visto ou cometido na prática (uma tabela que misturava
  níveis de detalhe e gerava um número errado que alguém levou para uma reunião).
  Nomes e valores só em ordem de grandeza.
}

A recompensa por esse cuidado aparece na hora de consultar. Com o grão claro e as
dimensões bem desenhadas, as perguntas de negócio viram consultas curtas e previsíveis:
junta o fato às dimensões de que a pergunta precisa, filtra, agrupa e soma.

\begin{sql}
-- Receita por mês e por categoria de produto, no esquema estrela de vendas
SELECT
    d.ano,
    d.mes,
    p.categoria,
    SUM(f.valor_total) AS receita
FROM fato_vendas  AS f
JOIN dim_data     AS d ON d.id_data    = f.id_data
JOIN dim_produto  AS p ON p.id_produto = f.id_produto
GROUP BY d.ano, d.mes, p.categoria
ORDER BY d.ano, d.mes, p.categoria;
\end{sql}

Essa previsibilidade também ajuda a IA. Quando o livro tratar de dados como contexto
para modelos de linguagem (capítulo~\ref{cap:contexto-ia}), o modelo dimensional bem
feito vai aparecer como meio caminho andado: grão explícito e nomes claros reduzem a
chance de a máquina somar a coisa errada.

\section{Inmon ou Kimball?}
\label{sec:inmon-ou-kimball}

Durante anos, a pergunta dividiu a área. A Tabela~\ref{tab:inmon-kimball} resume as
diferenças principais.

\begin{table}[htbp]
  \centering
  \caption{Inmon e Kimball lado a lado.}
  \label{tab:inmon-kimball}
  \small
  \begin{tabularx}{\textwidth}{@{}>{\raggedright\arraybackslash}p{0.2\textwidth}
      >{\raggedright\arraybackslash}X>{\raggedright\arraybackslash}X@{}}
    \toprule
    & Inmon & Kimball \\
    \midrule
    Ponto de partida & O armazém corporativo, de cima para baixo & Os processos de negócio, de baixo para cima \\
    \addlinespace
    Modelo do núcleo & Normalizado (3FN) & Dimensional (esquema estrela) \\
    \addlinespace
    Data marts & Dependentes, derivados do armazém & São a própria camada de apresentação \\
    \addlinespace
    Integração & No armazém central & Pelas dimensões conformadas (barramento) \\
    \addlinespace
    Primeira entrega & Mais lenta, pede o modelo corporativo antes & Mais rápida, um processo por vez \\
    \addlinespace
    Prioridade & Consistência e integração na empresa toda & Facilidade de consulta para o usuário \\
    \bottomrule
  \end{tabularx}
  \fonte{elaborado pelo autor com base em \textcite{inmon2005} e \textcite{kimball2013}.}
\end{table}

Na prática, muitas empresas acabaram com uma mistura das duas: uma camada integrada
perto das fontes, que lembra o armazém de Inmon, e modelos dimensionais para consumo,
no estilo de Kimball. Nos anos 2000 surgiu ainda uma terceira opção para a camada de
integração, o Data Vault, que separa as chaves de negócio, os relacionamentos e os
atributos descritivos em estruturas próprias, pensando em fontes que mudam com
frequência \parencite{linstedt2015}.

A nuvem mudou a conta de novo. Com armazenamento barato e processamento que cresce sob
demanda (capítulo~\ref{cap:nuvem}), algumas preocupações de desempenho que pesavam nas
decisões de Inmon e Kimball perderam força. É comum ver tabelas largas, que repetem
atributos que um modelo normalizado evitaria. Mesmo assim, a modelagem dimensional
continua no centro dos projetos de analytics engineering, porque o problema que ela
resolve não é de desempenho: é fazer o dado ser entendido do mesmo jeito por todo mundo
que o consulta.

\rascunho{
  Opinião do autor, com base no que viu funcionar: qual das abordagens aparece na
  prática dos projetos de hoje e por quê. Pode fechar com um relato curto.
}

\section{ETL, cargas noturnas e cubos OLAP}
\label{sec:etl-olap}

Os armazéns de dados dos anos 1990 e 2000 eram abastecidos por processos de ETL:
extrair os dados das fontes, transformá-los fora do banco de destino, num servidor
próprio, e carregá-los já prontos. Como as fontes eram sistemas operacionais que não
podiam ficar lentos durante o expediente, a carga acontecia de madrugada, numa janela
que precisava terminar antes de a empresa abrir. Quando a janela estourava, os
relatórios da manhã chegavam com o dado do dia anterior ao anterior.

Para que a consulta fosse rápida, os dados agregados eram pré-calculados em cubos
multidimensionais: estruturas em que cada combinação de dimensões (mês, região,
produto) já tinha o total guardado. O termo OLAP, de processamento analítico em linha,
foi proposto em 1993 num documento assinado por Edgar Codd, o criador do modelo
relacional, com doze regras para esse tipo de ferramenta \parencite{codd1993}. O
documento foi encomendado por um fornecedor de software de cubos, o que gerou
críticas sobre a independência das regras. O termo, ainda assim, pegou.

Em volta dos armazéns e dos cubos cresceram as suítes corporativas de BI, que
juntavam relatórios, cubos e painéis num mesmo pacote, com licenças caras e equipes
especializadas. No fim dos anos 2000, os maiores fornecedores independentes dessas
suítes foram comprados por grandes empresas de software corporativo.
\pendente{conferir e citar as aquisições de 2007 e 2008}

Desse período, o que envelheceu foi a infraestrutura: janelas noturnas, cubos
pré-calculados, servidores dedicados. O que ficou foi o modelo. Quando o processamento
ficou barato o bastante para calcular os totais na hora, os cubos perderam a razão de
existir, mas os fatos, as dimensões e o grão continuaram exatamente onde estavam. O
próximo capítulo mostra o que aconteceu quando a ferramenta de análise saiu do
departamento de TI e foi parar na mesa do analista.

\begin{leituras}
  \leitura{kimball2013}{A referência da modelagem dimensional. Os primeiros capítulos
    apresentam o processo de quatro passos; os demais aplicam a técnica a casos de
    vários setores.}
  \leitura{inmon2005}{A visão de Inmon na fonte: a definição de data warehouse, o modelo
    corporativo normalizado e o cuidado com a granularidade dos dados.}
  \leitura{inmon2001}{A Corporate Information Factory, a arquitetura completa em que o
    data warehouse é uma das peças.}
  \leitura{linstedt2015}{Para conhecer o Data Vault, alternativa para a camada de
    integração quando as fontes mudam com frequência.}
  \leitura{devlin1988}{Artigo curto e histórico, um dos primeiros a descrever a
    arquitetura de um armazém de dados separado dos sistemas operacionais.}
\end{leituras}
